\chapter{Making GDB Comfortable}\label{ch:gdb-comfort}

Out of the box GDB is a bare prompt with no source on screen and no memory of
your preferences. Twenty lines of configuration fix both, and are worth writing
once.

\section{The TUI}\label{sec:gdb-tui}

\begin{definition}[TUI]
	GDB's built-in Text User Interface: a split terminal with the source, the
	registers or the disassembly above and the command prompt below. It is part of
	GDB, needs nothing installed, and follows the program as it runs.
\end{definition}

\begin{keys}[0.30]
	\cmd{gdb -tui ./app}     & Start with it on                                     \\
	\key{<C-x>} \key{a}      & Toggle it at any time                                \\
	\key{<C-x>} \key{1}      & One window: source                                   \\
	\key{<C-x>} \key{2}      & Two windows: source and disassembly                  \\
	\cmd{layout src}         & Source only                                          \\
	\cmd{layout asm}         & Disassembly                                          \\
	\cmd{layout split}       & Both                                                 \\
	\cmd{layout regs}        & Add the registers                                    \\
	\key{<C-x>} \key{o}      & Move focus between windows                           \\
	\key{<C-l>}              & Redraw, when the display gets corrupted              \\
	\cmd{focus cmd}          & Give the arrow keys back to the command line         \\
	\cmd{tui reg general}    & Which register group to show                         \\
\end{keys}

\begin{note}
	The first thing everyone hits: with focus in the source window, the arrow keys
	scroll the source instead of walking the command history. \cmd{focus cmd}, or
	\key{<C-x>} \key{o}, gives them back. The second thing: program output and the
	TUI share a terminal and will corrupt each other --- \key{<C-l>} redraws, and
	\cmd{set inferior-tty /dev/pts/3} sends the program's output to a separate
	terminal for good.
\end{note}

\begin{moral}
	If you do nothing else in this chapter, learn \key{<C-x>} \key{a}. Seeing the
	current line highlighted in context, updating as you step, is the difference
	between GDB feeling archaic and feeling like a debugger.
\end{moral}

\section{\texorpdfstring{\file{.gdbinit} and User Commands}{.gdbinit and User Commands}}\label{sec:gdb-gdbinit}

GDB reads \file{\textasciitilde/.gdbinit} at startup, and then \file{./.gdbinit}
from the current directory if it is trusted.

\begin{gdbcode}
# ~/.gdbinit
set confirm off                  # do not ask "are you sure" for everything
set pagination off               # do not stop every screenful
set history save on              # remember commands between sessions
set history size 10000
set print pretty on              # structs one member per line
set print array on
set print elements 0             # never truncate
set disassembly-flavor intel     # if you prefer Intel syntax
set breakpoint pending on        # do not ask about future shared libraries
set backtrace limit 100          # tame runaway recursion

# a source window by default
tui enable
\end{gdbcode}

\begin{note}
	A project-local \file{.gdbinit} is read only if its directory appears in GDB's
	auto-load safe path, which is a security measure --- otherwise
	cloning a repository and running \cmd{gdb} in it would execute someone else's
	commands. Add \cmd{set auto-load safe-path /} to your global file if you trust
	everything you check out, and understand what you are turning off.
\end{note}

User-defined commands make a repeated inspection into one word:

\begin{gdbcode}
define plist
  set $n = $arg0
  while $n != 0
    printf "%s\n", $n->name
    set $n = $n->next
  end
end
document plist
Walk a linked list from the given head, printing each name.
end

define bpl
  break $arg0
  commands
    silent
    printf "hit %s\n", "$arg0"
    bt 3
    continue
  end
end
\end{gdbcode}

\begin{keys}[0.26]
	\cmd{\$arg0}, \cmd{\$arg1} & The command's arguments                            \\
	\cmd{\$argc}               & How many were given                                \\
	\cmd{document name}        & Text for \cmd{help name}                           \\
	\cmd{source debug.gdb}     & Load a file of commands mid-session                \\
	\cmd{define hook-run}      & Run automatically before every \cmd{run}           \\
\end{keys}

\begin{moral}
	Every project accumulates two or three inspections you do constantly --- print
	this tree, dump that ring buffer, check those invariants. Put them in a
	\file{.gdbinit} beside the source and commit it. It is documentation that
	executes.
\end{moral}

\section{Pretty-Printers and Python}\label{sec:gdb-python}

GDB embeds Python, and the most visible use is \vocab{pretty-printers}:
functions that decide how a type is displayed.

\begin{gdbcode}
(gdb) print v                                    # without a pretty-printer
$1 = {_M_impl = {<std::allocator<int>> = {...}, _M_start = 0x5555555592a0,
  _M_finish = 0x5555555592ac, _M_end_of_storage = 0x5555555592b0}}
(gdb) print v                                    # with one
$2 = std::vector of length 3, capacity 4 = {1, 2, 3}
\end{gdbcode}

\begin{note}
	For the C++ standard library these ship with GCC and are usually loaded
	automatically --- \cmd{info pretty-printer} lists what is active, and
	\cmd{set print pretty-printer off} turns them off when you need to see the raw
	struct. For your own types, a printer is a short Python class registered
	against a type name; \cmd{help python} and the GDB manual's \emph{Pretty
		Printing} chapter have the pattern.
\end{note}

\begin{keys}[0.38]
	\cmd{python print(gdb.parse\_and\_eval("n"))} & Read a program value from Python \\
	\cmd{python gdb.execute("bt")}                & Run a GDB command from Python    \\
	\cmd{info pretty-printer}                     & What is registered                \\
	\cmd{source printers.py}                      & Load a Python script              \\
\end{keys}

\section{Front Ends, and GDB from Neovim}\label{sec:gdb-frontends}

\begin{keys}[0.30]
	\cmd{gdb -tui}       & Built in, always available, no setup                    \\
	\cmd{cgdb}           & A curses front end with a better source window          \\
	\cmd{gdbgui}         & A browser front end                                     \\
	\texttt{nvim-dap}    & Debug Adapter Protocol client for Neovim                \\
	\texttt{termdebug}   & Ships \emph{with} Vim and Neovim --- \cmd{:packadd termdebug} \\
\end{keys}

\begin{eg}[\texttt{termdebug}]
	The lowest-effort way to debug inside the editor of Part~\ref{part:vim}: it is
	already installed, and needs no plugin.
\begin{vimcode}
:packadd termdebug
:Termdebug ./app
\end{vimcode}
	Neovim splits into a GDB window, a program-output window and your source. Then
	\cmd{:Break} sets a breakpoint on the current line, \cmd{:Run}, \cmd{:Step},
	\cmd{:Over} and \cmd{:Finish} drive it, and \cmd{:Evaluate} prints the
	expression under the cursor --- while the source buffer stays an ordinary
	Neovim buffer with all the motions of \autoref{ch:vim-motions}.
\end{eg}

\begin{note}
	\texttt{nvim-dap} is the fuller answer --- a proper debug adapter with a
	variables pane, breakpoint signs in the gutter and per-language configuration
	--- and it is a plugin to add through \texttt{lazy.nvim}
	(\autoref{sec:nvim-plugins}). It is not installed on this machine. Start with
	\texttt{termdebug}, which costs nothing, and add \texttt{nvim-dap} if you find
	yourself debugging daily.
\end{note}

\section{A Worked Debugging Session}\label{sec:gdb-session}

Everything in Part~\ref{part:gdb}, on one bug. The program counts words in a
file and reports zero for a file that plainly has some.

\begin{ccode}
int count_words(const char *path) {
    FILE *f = fopen(path, "r");
    char buf[64];
    int n = 0;
    while (fgets(buf, sizeof buf, f)) {
        for (char *p = strtok(buf, " \n"); p; p = strtok(NULL, " \n"))
            n++;
    }
    fclose(f);
    return n;
}
\end{ccode}

\begin{gdbcode}
$ gcc -g -O0 -o wc main.c && gdb -q ./wc
(gdb) break count_words
(gdb) run input.txt
Breakpoint 1, count_words (path=0x7fffffffe5c9 "input.txt") at main.c:6
6           FILE *f = fopen(path, "r");
(gdb) next
8           int n = 0;
(gdb) print f
$1 = (FILE *) 0x5555555592a0          # the file opened fine
(gdb) display n                       # watch the counter as we go
1: n = 0
(gdb) break 10 if n > 0               # stop as soon as it counts anything
(gdb) continue
Continuing.
1: n = 0                              # ...it never does
[Inferior 1 (process 41233) exited normally]
\end{gdbcode}

The loop body never runs, so \cmd{fgets} returned \texttt{NULL} the first time.
Second run, with the attention one line earlier:

\begin{gdbcode}
(gdb) run input.txt
(gdb) break 9
(gdb) continue
Breakpoint 2, count_words (...) at main.c:9
9           while (fgets(buf, sizeof buf, f)) {
(gdb) print sizeof buf
$2 = 64
(gdb) print fgets(buf, 64, f)
$3 = 0x0                              # NULL on the very first read
(gdb) print ftell(f)
$4 = 0                                # ...yet we are at the start of the file
(gdb) print feof(f)
$5 = 0
(gdb) print ferror(f)
$6 = 1                                # the stream is in an error state
\end{gdbcode}

\begin{note}
	\cmd{print fgets(...)} and \cmd{print ferror(f)} are \emph{calling} those
	functions in the live process (\autoref{sec:gdb-force}). Being able to
	interrogate a C library from the prompt --- with the program's own \texttt{FILE
		*} in hand --- is what makes GDB more than a viewer.
\end{note}

The stream errored without reaching end of file, which points at the mode.
\cmd{print path} shows \texttt{"input.txt"}, and a look at the caller resolves it:
the caller opened the same file earlier, closed it, and passed the stale name
while an \cmd{fopen} failure elsewhere left \texttt{errno} set. The fix is the
check that was never there:

\begin{ccode}
FILE *f = fopen(path, "r");
if (!f) { perror(path); return -1; }
\end{ccode}

\begin{moral}
	Note what the session did \emph{not} involve: no recompiling with print
	statements, no guessing, and no reading of the whole program. Establish where
	it stops being right, look at the state there, and walk outwards one frame at a
	time. That loop is the entire method, and the commands in
	Part~\ref{part:gdb} are only its vocabulary.
\end{moral}

\begin{tryit}
	Write a ten-line C program with a deliberate null dereference. Compile it with
	\cmd{gcc -g -O0}, run \cmd{gdb -q ./a.out}, then \cmd{run}, \cmd{bt},
	\cmd{info locals}. You have now done the thing that most people avoid for
	years, and it took under a minute.
\end{tryit}
